\begin{abstract}

UMI develops ASL-to-English translation through two complementary forms of
competition. Endpoint miners serve translations using systems they control.
Model contributors submit reproducible, licensed model bundles that UMI can
preserve and operate independently. Both are evaluated against held-out ASL
examples. Verified improvements can become the next public reference model.

The network must produce useful translations and an accumulating body of
runnable models. Access to a successful endpoint alone cannot meet the second
objective. Model promotion therefore requires possession of the model files,
a reproducible inference environment, documented rights, and verification
without the original miner or an external inference API.

The first successor release separates implementation rehearsal from reward
activation. No existing miner must disclose proprietary weights retroactively.
The private endpoint track remains available under its published rules; the
model-contribution track has explicit publication and reuse terms.

\end{abstract}

\section{Scope and transition}\label{1-scope-and-transition}

The task is raw American Sign Language video to bounded English text.
A miner may use an end-to-end video model, pose model, ensemble or hosted API
for endpoint service, subject to the service deadline and applicable terms.
A promoted model must run with its archived dependencies and no external
inference service. Future tasks require their own evaluation and reward policy.

The frozen two-miner pilot bootstrap is retired. Its signed manifest and
historical evidence remain unchanged; see the
\href{https://github.com/Umi-BitSign/umi/blob/main/docs/BOOTSTRAP_WEIGHT_ADDENDUM.md}{bootstrap addendum} and
\href{https://github.com/Umi-BitSign/umi/blob/main/docs/SHARED_VALIDATOR_BOOTSTRAP_SUPERSESSION.md}{shared-validator supersession}.
It was replaced by the separately signed
\href{https://github.com/Umi-BitSign/umi/blob/main/docs/REGISTRATION_BRIDGE.md}{temporary registration bridge}.
The \href{https://github.com/Umi-BitSign/umi/blob/main/docs/REGISTRATION_BRIDGE_FUNDING_CAP_2026-09-13.md}{September 13 deployment record}
records funding-grouped bridge rows finalized by validators UID 0 and UID 54.

The bridge checks registered miners\textquotesingle{} chain-announced HTTPS health endpoints.
Prior pilot participation and a running translation model are not required.
Passing miners connected by a shared coldkey, HTTPS endpoint IP or recorded
pre-registration funding source share one equal group weight budget, divided
among their passing UIDs. Validator-permitted and owner-associated identities
remain excluded. The \href{https://github.com/Umi-BitSign/umi/blob/main/docs/CURRENT_MINER_OPERATION.md}{miner guide} gives the
complete eligibility and endpoint requirements.

Funding assertions bind an exact registration identity in the signed policy.
A cached coordinator checker detects new registrations; new funding links
affect weights only after a signed snapshot refresh. Validators and miners
do not each need a Taostats API key. Shared-sender grouping can combine
independent exchange customers, and separate funding sources can evade it.
It does not prove common human ownership. The
\href{https://github.com/Umi-BitSign/umi/blob/main/docs/REGISTRATION_BRIDGE_FUNDING_CAP.md}{funding-cap rule} is temporary and
is not a requirement of the 70/30 successor mechanism specified below.

Bridge weights measure availability only. They establish neither translation
quality nor successor qualification. A separately signed version3 bridge policy
can continue until explicitly superseded, with no scheduled expiry. Historical
version1/2 policies retain their original cutoffs. A new lifetime authorization
and compatible release are required; this document alone does not activate an
extension, reopen the retired pilot or launch translation competition.

Version 0.2 supersedes \href{https://github.com/Umi-BitSign/umi/blob/main/whitepaper/LEGACY_V0_1.md}{version 0.1}\textquotesingle s proposed fixed four-validator,
three-publisher and 30-day-soak launch gates. It replaces them with a published
evaluation cohort, disclosed control groups, bounded rehearsals, policy-selected
result reproduction and the explicit activation checklist in Section 10.
This change is a UMI policy revision. It is not evidence that any earlier
activation gate was completed.

The active mechanism and policy hash must be distinguishable from proposed
rules in the website, observer API, miner guide and signed validator feed.
Software installation, a successful test, and positive native chain emissions
are not interchangeable with activation of this competition.

\section{Participation and identity}\label{2-participation-and-identity}

A publisher prepares and publishes challenge batches under the data policy.
A validator\textquotesingle s supervisor is the host service that checks signed release
directives and manages its worker. The worker image is the versioned container
containing the validator task code.

Anyone with a registered SN78 hotkey can submit through the same public
interface. A prior pilot, invitation, GitHub issue or private upload key is
not an admission requirement. Registration proves control of an identity;
it does not establish translation quality or a right to rewards.

A submission binds:

\begin{itemize}
\tightlist
\item
  network, subnet, competition policy hash and track;
\item
  the submitting hotkey and a monotonic submission sequence;
\item
  the model or endpoint revision and applicable publication terms;
\item
  the validity interval in finalized blocks; and
\item
  the complete artifact manifest for a model contribution.
\end{itemize}

The hotkey signs a domain-separated digest of canonical JSON. Signatures from
another network, policy, identity or track cannot be reused. The registry
returns the same receipt for an identical retry. A conflicting sequence is
rejected. The policy sets the minimum interval between accepted replacements
and admits at most one current submission per hotkey per track.

The registry publishes accepted submissions and explicit rejection reasons.
Pending work survives process restarts. A UID is resolved from the current
finalized hotkey mapping at the scoring boundary and rechecked before a
weight submission; UID reuse must not transfer another miner\textquotesingle s score.
Deregistration, permit changes and revision replacements cannot rewrite a
closed evaluation roster.

The coordinator must publish the complete admission log and each round\textquotesingle s
roster. Queue receipt time, accepted block and expiry remain distinct.
A delayed scheduler cannot backdate a usable evaluation window or count
an unissued request as a miner failure. Closed rounds are immutable.
Capacity limits and the deterministic selection rule must be public.

\section{Two tracks}\label{3-two-tracks}

\subsection{Endpoint service}\label{31-endpoint-service}

An endpoint miner runs the UMI translation protocol using its own model or
provider. The miner retains its weights and implementation unless it separately
opts into model contribution. A health-only keepalive earns no translation score.

Each accepted endpoint revision is challenged within a published request
budget. Requests and answers are authenticated and bound to the round, clip,
policy and revision. Raw video is verified before use. A miner timeout or
invalid response scores zero for that assigned case. Coordinator, publisher
or evaluator infrastructure failure voids the affected evaluation instead of
being charged to the miner.

An unavailable miner does not stall the entire cohort. The next complete
round calculates its row from valid evidence and public failure rules.
Retries cannot replace an observed answer with a more favorable answer.

\subsection{Model contribution}\label{32-model-contribution}

A model contributor publishes a versioned bundle under the policy\textquotesingle s accepted
reuse terms. Its content identity is independent of the submitting UID,
repository name and upload time. The same bytes submitted under another
hotkey do not become a new model.

A complete contribution contains:

\begin{itemize}
\tightlist
\item
  weights, including the exact base weights needed by an adapter or LoRA;
\item
  architecture and inference configuration;
\item
  tokenizer, vocabulary, processor and all preprocessing assets;
\item
  inference code and a pinned environment;
\item
  license texts, attribution, upstream dependency identities and provenance;
\item
  the parent public baseline identity, if derived from it; and
\item
  instructions for a clean, offline inference run.
\end{itemize}

A repository URL or a hash-only declaration is insufficient for promotion.
The archive verifier reads every declared byte and rejects missing, extra,
oversized or changed files, path traversal, links and special files.
External paid inference and a dependency on the original miner\textquotesingle s server
disqualify the bundle from the offline promotion profile.

Automatic admission validates structure, signatures and declared terms.
It does not prove authorship, consent or compatibility of upstream licenses.
Promotion includes review of those rights. Review happens on the submitted
bundle; enrollment must not require private credentials or wallet material.

\section{Evaluation}\label{4-evaluation}

\subsection{Data and holdouts}\label{41-data-and-holdouts}

The launch uses automatic scoring against a private, labeled ASL holdout. It
does not require a person to grade miner outputs or a new ASL review panel before
activation. Existing dataset annotations may supply the English references when
their source, annotation procedure and permitted use are documented. Dataset
validation and model-contribution rights review remain required; neither is a
per-output human scoring step.

Reserve the evaluation split before training or tuning. Clips, labels and
near-duplicate recordings used to train or select the baseline or a candidate
cannot count as unseen evaluation evidence for that model. Making training data
private afterward does not make it a holdout. Record known training overlap and
unknown upstream exposure rather than claiming they have been ruled out.
Publish a development set separately. Hidden evaluation clips and labels must
not become miner training data through repeated tests, debug logs, image layers
or unrestricted evaluator APIs.

A round commits its exact policy, candidate roster, incumbent baseline,
evaluation suite and runtime profile. Candidates and incumbent are scored
against the same suite. A score from a different suite or baseline cannot
justify promotion. Suites are retired after their protected evaluation
interval and replaced before adaptive resubmissions can exploit them.

Evaluation evidence binds video hashes, accepted references, outputs,
failures, timing, environment and evaluator identities. The public replay
package reveals only material whose publication is permitted by its consent
and data policy. When underlying clips cannot be redistributed, disclose that
limitation and retain access-controlled independent verification; do not call
a hash alone public reproducibility.

Data collection and labeling can be developed separately. They receive no
implicit share of this policy\textquotesingle s miner rewards. The
\href{https://github.com/Umi-BitSign/umi/blob/main/docs/OPEN_COMPETITION_PRIVATE_HOLDOUT.md}{private-holdout launch procedure}
describes required inputs, annotation checks, split isolation and suite retirement.

\subsection{Deterministic quality}\label{42-deterministic-quality}

Reuse UMI\textquotesingle s exact text normalization and CER/WER implementation. The approved
initial launch uses \path|umi-open-competition-policy/2| with
\path|umi-competition-suite/2|: each case has exactly one authentic committed English
reference. Historical version 1 keeps its three-to-five-reference contract.
CER measures edit distance between normalized graphemes, excluding whitespace;
WER uses normalized word tokens. For each reference, let \path|d| be the edit
distance and \path|n| its number of scoring units. Its similarity is
\path|max(0, 1 - d / max(1, n))|. The case score is the highest similarity across
the committed references.

The initial version 2 profile uses CER for fingerspelling with exact weight
\path|3/13|, and WER for continuous signing with exact weight \path|10/13|. These preserve
the original 15:50 relative weights across the two available tasks. Short
utterances are outside this launch profile. Version 1 retains its 15\%/35\%/50\%
three-task weights. Compute the arithmetic mean of the case scores within each
stratum, then sum those means multiplied by their respective weights.
Every required stratum must meet the policy\textquotesingle s minimum case count.
Suites contain at least three cases in total. Policy and suite versions must
match; unknown or extra strata are rejected. The two-task score is not directly
comparable to a historical three-task score. This change does not alter the
70/30 reward split or the approved 120-second inference limit.
The normalization and scoring functions are defined in
\href{https://github.com/Umi-BitSign/umi/blob/main/src/umi/scoring.py}{scoring.py}.

Record integer case errors and exact rational aggregates. Floating-point
rounding cannot determine a promotion or a tie. The policy declares:

\begin{itemize}
\tightlist
\item
  minimum qualifying quality;
\item
  the absolute improvement margin for promotion;
\item
  the inference time and output-size limits;
\item
  the number of cases per stratum and evaluator control-group quorum; and
\item
  submission rate, artifact-size and queue bounds.
\end{itemize}

A promoted candidate must improve the aggregate by at least the declared
margin and must not regress in any required stratum. It must pass the offline
resource and reconstruction checks as well. A tie retains the incumbent.
No score is a claim of clinical fitness or a replacement for product-specific
safety evaluation. CER/WER measure agreement with the committed English
references. Without separate semantic validation, report these as automatic
benchmark scores, not human-verified ASL understanding. Valid paraphrases and
meaning-changing errors can be misranked by text edit distance.

\subsection{Evaluators and trust}\label{43-evaluators-and-trust}

A policy lists evaluator hotkeys and their actual administrative control groups.
Evaluators are automated execution and signing services; their operators do not
need to grade ASL outputs. Multiple hotkeys controlled by one operator count as
one group. UID 0 and UID 54 under the same administration cannot supply two
independent votes.

The approved initial launch selects only UID 0, hotkey
\path|5Fk765B4CRBekwErwE5VxvveWhHztHSfsnsLt8cbDayDWsuk|, in one disclosed UMI-operated
control group, with \path|required_evaluator_groups = 1|. This initial phase trusts
one operator\textquotesingle s automated evaluation and publication; it does not provide
independent cross-operator reproduction or resistance to that operator\textquotesingle s
collusion. UID 54 is not a second evaluator or vote. Adding evaluators or raising
the quorum requires a later explicit signed policy and deployment rehearsal.
See the \href{https://github.com/Umi-BitSign/umi/blob/main/docs/OPEN_COMPETITION_UID0_LAUNCH.md}{UID 0 launch profile}.

Evaluators sign the same canonical evaluation result. A quorum certificate
contains valid signatures from the policy\textquotesingle s required number of distinct
control groups. Multiple keys from one group cannot increase that count.
A small disclosed cohort is a trust assumption, not permissionless evaluation.

The current rehearsal schema represents one common recorded transcript:
candidate and incumbent outputs, outcome statuses, elapsed milliseconds,
completion block and the pinned runtime identity. All of those fields are
inside the signed result. A co-signature endorses that transcript; it does not
claim that the signer measured identical wall-clock timing in a separate run.
The result object remains unchanged. A separate evidence envelope now binds
one independently signed run record to each common-result signer. Each record
binds the policy, round, submission, common result, suite, model and runtime,
with its own interval, timings and execution-evidence digest. Outputs and
per-case resource eligibility must agree with the common transcript, and exact
score replay must agree. The run signer set must match the common certificate
with one key per authorized control group. This contract has local tests;
production execution and evidence retention still require rehearsal. In the
initial single-operator phase there is one signed run record, not an independent
second run. Later multi-group policies require each selected group to reproduce
the result.
Co-signatures and execution-evidence digests alone do not prove a run occurred.

The local paired-model runner now records bounded raw outputs and finalized
boundary references in a separate wallet-free execution journal. It reserves an
attempt before running model code, retains returned outputs before the next
finality read and refuses automatic reruns after interruption. Completed retries
return the original evidence. Local aggregation checks the revealed suite and
distinct evaluator groups, proposes the maximum observed timing per case, and
prepares unsigned run records. This does not supply production round
publication, independent operators or authorization to submit weights.

The evidence identity is the policy, round and miner-submission digest.
Store each authenticated quorum result and its signatures durably. Reordering
signatures or adding signatures for the same result is an idempotent retry.
Two different quorum-certified results for one evidence identity put the
entire round on conflict hold, even when the quorums have no group in common.
Preserve both results and any proof that a control group signed both.
Another hotkey in that group cannot erase its prior statement.

Invalid signatures, incorrect bindings and a statement without the required
quorum cannot trigger this automatic hold. Ordinary non-quorum disagreement
does not give a minority a veto. Any separate group-exclusion rule requires
an explicit policy; it cannot be invented while settling a disputed round.

Authentication checks the historical assignment and result interval, separate
from present eligibility or whether the result earns a score. A valid failed
outcome, expired reward window or failed archive check cannot erase evidence
of conflicting statements. Commit that evidence before reporting a rejected
action. Promotion and reward projection read the recorded results and check
the round\textquotesingle s conflict state within their action transaction, including retries.

Missing quorum evidence prevents settlement. Before activation, the policy
must fix the evidence-collection cutoff and the procedure for publishing an
immutable settlement record. A local projection cannot establish that all
evidence has arrived. The local store fixes a cutoff before round close,
retains first-observation blocks and settles only evidence recorded by that
cutoff. Its immutable artifact binds the roster, results, run evidence, suite,
registration snapshot, promotion head and projected row. Local block inputs
and records are not finalized observations or signed publication authority.
Late conflicting certificates remain admissible as
incident evidence and stop further use of the disputed round. They cannot
rewrite historical records or undo payments already finalized on chain.

Retain requests and outputs so a verifier can recompute scores. Inference
reproduction is distinct from arithmetic replay: hardware-specific numerical
variation must be investigated under a pinned runtime profile rather than
hidden behind approximate score agreement.

\subsection{Untrusted execution}\label{44-untrusted-execution}

Downloaded artifacts never execute inside a wallet-bearing validator process.
Use a separate, resource-limited evaluation environment with no wallet mounts,
host sockets, private keys, cloud credentials or plaintext reference labels.
Allow only the assigned public inputs and a bounded output channel.

The offline profile disables network access, uses read-only model files,
pins the evaluator runtime, and enforces process, CPU, memory, storage,
output and wall-clock limits. A timeout terminates the complete workload.
Merely starting a container does not prove these controls are effective;
release tests must demonstrate them on the supported host.

Large VLMs may require GPU evaluators. Existing CPU-only validator hosts are
not promised the ability to run arbitrary contributed models. Each activated
policy pins a supported evaluation budget. Training hardware remains the
miner\textquotesingle s choice.

\section{Preservation and baseline promotion}\label{5-preservation-and-baseline-promotion}

A baseline registry is an append-only sequence of promotions. Each entry binds
the previous baseline, new model manifest, policy, evaluation result, agreeing
evaluator signatures, archive verification and rights review. The initial
baseline is imported explicitly with its own manifest and provenance.
The imported initial baseline has no registered contributor and earns no
model-contribution share. Until the first qualifying promotion, a policy must
allocate zero basis points to that track or no successor reward row can be
produced.

Promotion follows this order:

\begin{enumerate}
\def\labelenumi{\arabic{enumi}.}
\tightlist
\item
  Verify the submitted identity and complete candidate bundle.
\item
  Reconstruct and run the model offline within the declared resource budget.
\item
  Replay the paired evaluation and verify evaluator-group agreement.
\item
  Verify rights and archive the exact files in UMI-controlled storage.
\item
  Commit a compare-and-swap update against the evaluated incumbent.
\item
  Publish the new baseline and its promotion evidence.
\end{enumerate}

If another promotion changed the incumbent, the candidate needs a new paired
evaluation. An old favorable result cannot overwrite a newer baseline.
An interrupted archive or failed verification leaves the prior baseline
unchanged. Repeating the same successful promotion is idempotent.
That retry remains subject to the conflict guard. If late evidence disputes a
completed promotion, preserve its record and archived files, but hold further
promotion and reward projection using that baseline or its descendants.
Resuming requires an explicit recovery transition; a newer round alone cannot
clear the hold or silently reassign the original contribution.

A promoted version remains retrievable when the miner disappears or changes
its repository. Maintain a separate backup and periodically test restoration.
Hugging Face is a distribution channel; a mutable branch or a single external
repository is not the preservation record.

Future miners may download and improve the current baseline. They must retain
required attribution and identify their parent version. Cosmetic repackaging
or a copied checkpoint cannot earn a fresh promotion reward. More general
model-copy detection is imperfect and must not be presented as cryptographic
proof of independent training.

\section{Licenses, ownership and product use}\label{6-licenses-ownership-and-product-use}

Publication grants only the rights supplied by the applicable licenses.
It does not assign exclusive intellectual-property ownership to UMI.
Require rights sufficient for preservation, commercial inference,
modification and redistribution, with compatible upstream obligations.
Exclusive acquisition, when desired, needs a separate agreement.

The current reference repository labels its code Apache-2.0 and its weights
and portable bundle CC BY-SA 4.0. Those obligations are not erased when a miner
fine-tunes, changes hosting, or submits a model. Preserve license text,
attribution and provenance with every promoted version. Do not silently
relabel inherited weights under a different license.

A contributor must identify base-model and dataset sources and attest that
it can grant the proposed rights. UMI must review material restrictions
before promotion. Dataset consent, privacy and third-party restrictions
remain separate from permission to run inference.

Before contribution intake opens, publish the accepted-license list, exact
contribution terms, review route and required provenance evidence. The signed
policy binds the list and terms; matching a license identifier is not rights
clearance. The \href{https://github.com/Umi-BitSign/umi/blob/main/docs/MODEL_CONTRIBUTION_REVIEW.md}{contributor review checklist}
describes source lineage, permissions, reconstruction and decision records.
Research-only, non-commercial or unspecified source permissions need explicit
review; they are not automatically cleared by a new license on the bundle.
Unresolved material restrictions hold promotion.

Preliminary review should let contributors resolve source questions before
training. Record the disclosed versions, intended uses, policy, terms and any
conditions. Reuse unchanged reviewed evidence at final review, while allowing
reassessment for new facts or restrictions. Preliminary feedback does not
guarantee promotion or replace review of the submitted artifact. Endpoint
miners need not disclose private weights solely to serve requests, but remain
responsible for applicable model, data and service-use restrictions.

Public baselines can support decentralized replication and independently
operated products. Neither participation in Bittensor nor this endpoint
track requires every miner to surrender private weights.

\section{Reward calculation}\label{7-reward-calculation}

Reward allocation is a required signed-policy value, expressed as integer
basis points. There is no default allocation and this whitepaper does not
activate a split. Endpoint and model-contribution shares sum to 10,000.

The approved initial launch allocation is 7,000 basis points (70\%) for endpoint
service and 3,000 (30\%) for the current promoted model\textquotesingle s contributor. Both
reward tracks launch together. The signed activation policy must carry that
allocation, and activation requires a qualifying preserved promotion with an
eligible contributor. Approval of the allocation does not activate rewards or
change the historical bootstrap.

For each closed round, qualifying endpoint submissions receive their share
in proportion to their exact quality scores. Scores below the policy floor
and invalid miner responses receive zero. A model-contribution share is paid
to the registered contributor of the currently valid promoted baseline,
provided it has qualifying fresh evaluation evidence in that round.
Copying the public baseline does not transfer that attribution.
The evaluator can run the preserved incumbent directly. Its contributor does
not need to keep an inference server online or renew an old submission to
make those preserved bytes available; current registration is still required.

This is king-of-the-hill attribution for the model contribution share,
not winner-takes-all subnet emissions. A contributor may separately operate
an endpoint and qualify for the service share. There are no automatic
perpetual royalties after replacement.

If a nonzero track allocation has no qualifying recipient, the round cannot
produce a successor row under the initial profile. It must report the reason
and refrain from submission. It must not quietly donate the missing share,
normalize it onto another track, fabricate scores or redirect it to an
owner UID. A later policy may define a different explicit fallback.

Resolve recipients against a common finalized SN78 snapshot. Merge both
track contributions by current UID, normalize the complete vector using
exact rational arithmetic, and produce canonical u16 raw weights.
Publish both the rational allocation and resulting raw row. Validators
must recheck policy validity, registration and chain requirements before
signing. Finalized consensus and incentive determine actual payouts.

A projected row is not a submitted or economically effective row.
Observer status reports those states separately. The mechanism does not
guarantee a particular amount of income.

\section{Policy and state boundaries}\label{8-policy-and-state-boundaries}

All signed protocol objects use canonical JSON and domain-separated digests.
A changed reward split, evaluation suite policy, resource profile,
license policy or admission rule has a different policy identity.

Policy objects have a version, activation interval, predecessor identity
and explicit track allocation. No implementation reads an arbitrary chain
value into the policy and treats that as authorization. Chain requirements
are checked against the approved policy.

Public state consists of submissions, acceptance receipts, closed round
rosters, evaluation results, promotion records and projected or finalized
weight evidence. Private state consists of wallets, sealed holdouts and
operational credentials. Public status must not expose private state.

The registry and promotion store use durable transactions. A process crash
cannot partially replace the current baseline or consume a valid retry.
Evidence objects are content-addressed and immutable. The baseline pointer
uses compare-and-swap against its recorded predecessor.
Conflict evidence commits separately from the action it prevents, so an action
rollback cannot remove the hold. Public status distinguishes historical
promotion records from a baseline that is currently held for conflict.

Open submission is not open write access to the baseline registry.
Only a verified promotion transaction may advance the baseline.

\section{Threats and limitations}\label{9-threats-and-limitations}

The initial successor must address:

\begin{itemize}
\tightlist
\item
  Sybil submissions and copied public models, through bounded intake,
  content identity and original promotion attribution;
\item
  benchmark leakage and adaptive overfitting, through protected suites,
  fixed rosters, fresh paired evaluation and rate limits;
\item
  malicious model files and dependencies, through byte verification,
  execution isolation and offline reconstruction;
\item
  evaluator collusion, through disclosed control groups, independent
  agreement and replayable evidence;
\item
  storage loss, through verified complete archives and restore tests;
\item
  scheduler delay, through acceptance receipts, usable deadlines and
  explicit infrastructure failures; and
\item
  policy or update-key compromise, through signed bounded releases,
  independent operator verification and a documented emergency hold.
\end{itemize}

These controls have limits. A signed license declaration can be false.
Different hashes can represent copied behavior. A coordinator can censor
intake. Evaluators can collude. Audit records expose many failures but do
not make those trust assumptions disappear.

\section{Release and activation}\label{10-release-and-activation}

The implementation must first run without any weight-submit capability.
A complete rehearsal exercises a new miner submission, replacement,
paired evaluation, failure handling, archive restoration, promotion,
baseline download by another miner, and exact row replay after restart.

Before activating successor rewards, publish:

\begin{itemize}
\tightlist
\item
  the signed policy, explicit reward allocation and activation interval;
\item
  supported miner instructions and the complete admission interface;
\item
  the initial preserved baseline and accepted license profile;
\item
  the exact evaluator cohort, control-group disclosures and resource profile;
\item
  evaluation-data rights and annotation provenance, dataset validation, split
  isolation and protected-suite procedure, without a mandatory human grading panel;
\item
  successful endpoint and contributed-model end-to-end evidence, including
  signed run records from the selected evaluator cohort bound to the common result, and adversarial,
  timeout, restart, replay and failed-promotion cases;
\item
  the deterministic row calculation and its finalized-chain preflight;
\item
  the compatible signed validator release and tested upgrade path;
\item
  the evidence cutoff, settlement record and late-conflict recovery rules; and
\item
  monitoring, expiry, incident handling and rollback procedures.
\end{itemize}

The policy fixes the rehearsal duration and minimum evaluator evidence
before it begins. Rehearsal output cannot be relabeled as production evidence.
A change affecting scoring or trust invalidates the affected evidence and
requires a new bounded rehearsal.

The current supervisor validates bootstrap input schemas on the host.
Changing a worker image alone does not authorize a new manifest or profile.
Test the actual installed supervisor against the successor. Publish a host
upgrade only if required, with explicit operator consent and rollback.
Never run competing weight writers for the same hotkey.

Keep the UID 0 and UID 54 registration-bridge processes operating within their
signed policy until the approved transition or their submission deadline.
A signed transition coordinates the stop of each old worker and activation of
its successor, preserving authorization history, keys and transaction journals.
The bridge\textquotesingle s inherited hard sunset remains effective even if the successor is
delayed.

\section{Implementation status}\label{11-implementation-status}

This specification defines the successor target. It does not claim that
open enrollment, arbitrary model evaluation, promotion or successor rewards
are deployed. The repository\textquotesingle s
\href{https://github.com/Umi-BitSign/umi/blob/main/docs/OPEN_COMPETITION.md}{implementation checklist} distinguishes tested
code, local rehearsal tooling and remaining production integration.

The earlier transport, canonicalization, deterministic scoring, finality
and supervisor components can be reused where their contracts match.
Historical pilot success establishes neither a model contribution nor
the successor\textquotesingle s quality qualification.

\section*{References}\label{references}
\addcontentsline{toc}{section}{References}

\begin{itemize}
\tightlist
\item
  \href{https://github.com/Umi-BitSign/umi-reference-model}{UMI reference model}
\item
  \href{https://docs.oroagents.com/docs/miners/evaluation-lifecycle}{ORO evaluation lifecycle}
\item
  \href{https://github.com/NousResearch/finetuning-subnet/blob/master/docs/miner.md}{Nous model contribution workflow}
\item
  \href{https://creativecommons.org/licenses/by-sa/4.0/}{CC BY-SA 4.0 terms}
\end{itemize}
